\section{Ακρίβεια και Σύγκριση Μοντέλων Υποκατάστασης}
\label{sec:eval_surrogate}

\subsection{Ακρίβεια Πρόβλεψης και Διαγράμματα Ισότητας (Parity Plots)}
Η αξιοπιστία του επιπέδου Tier 2 εδράζεται στην ικανότητα του μοντέλου υποκατάστασης (surrogate) να προβλέπει με υψηλή ακρίβεια την αεροδυναμική απόδοση $L/D$ των υποψήφιων διαμορφώσεων χωρίς να εισάγει συστηματικό σφάλμα (bias) στην εξελικτική πορεία.

Στο Σχήμα~\ref{fig:surrogate_pred_vs_true} παρουσιάζεται το διάγραμμα ισότητας (parity plot) μεταξύ των πραγματικών τιμών καταλληλότητας από τον επιλυτή AeroSandbox VLM και των προβλέψεων του εκπαιδευμένου μοντέλου surrogate, συνοδευόμενο από την κατανομή των σφαλμάτων υπολοίπων (residuals).

\begin{figure}[htbp]
    \centering
    \includegraphics[width=0.92\textwidth]{figures/fig2_surrogate_pred_vs_true.png}
    \caption{Διάγραμμα ισότητας (Parity Plot) πραγματικής έναντι προβλεπόμενης καταλληλότητας $L/D$: (α) Συσχέτιση επίλυσης AeroSandbox VLM και πρόβλεψης Surrogate επί της ιδανικής διαγωνίου $y=x$, και (β) ιστόγραμμα σφαλμάτων υπολοίπων (residuals) με σχεδόν κανονική κατανομή γύρω από το μηδέν ($\mu_{\text{err}} \approx 0$).}
    \label{fig:surrogate_pred_vs_true}
\end{figure}

Η ισχυρή γραμμική συσχέτιση των σημείων επί της διαγωνίου $y=x$ και η συμμετρική, συγκεντρωμένη κατανομή των υπολοίπων επιβεβαιώνουν ότι το surrogate μοντέλο δεν υποεκτιμά ούτε υπερεκτιμά συστηματικά την αεροδυναμική ποιότητα των διαμορφώσεων.

\subsection{Σύγκλιση Σφάλματος Γενίκευσης με το Μέγεθος Δειγμάτων}
Για να διαπιστωθεί ο ελάχιστος απαιτούμενος αριθμός δειγμάτων προθέρμανσης (warmup samples) για την ενεργοποίηση αξιόπιστων προβλέψεων, μελετήθηκε η εξέλιξη του συντελεστή προσδιορισμού $R^2$ και του μέσου τετραγωνικού σφάλματος (MSE) ως συνάρτηση του πλήθους δειγμάτων εκπαίδευσης $N_{\text{train}} \in [30, 300]$.

Στο Σχήμα~\ref{fig:surrogate_error_convergence} και στον Πίνακα~\ref{tab:surrogate_learning_curve} απεικονίζεται η καμπύλη μάθησης των μοντέλων Gaussian Process (GP), Multi-Layer Perceptron (MLP) και Random Forest (RF).

\begin{figure}[htbp]
    \centering
    \includegraphics[width=0.92\textwidth]{figures/fig2_surrogate_error_convergence.png}
    \caption{Σύγκλιση σφάλματος γενίκευσης και καμπύλες μάθησης: Εξέλιξη του συντελεστή $R^2$ και του MSE σε συνάρτηση με τον αριθμό δειγμάτων εκπαίδευσης $N_{\text{train}}$.}
    \label{fig:surrogate_error_convergence}
\end{figure}

../../evaluation/figures/table_fig2_surrogate_error_convergence.tex

Όπως προκύπτει από τον Πίνακα~\ref{tab:surrogate_learning_curve}, η διαδικασία Gaussian Process παρουσιάζει εξαιρετική αποτελεσματικότητα δειγματοληψίας (sample efficiency), επιτυγχάνοντας $R^2 = 0.877$ με μόλις 100 δείγματα εκπαίδευσης και κορυφώνοντας στο $R^2 = 0.914$ στα 200 δείγματα. Αντίθετα, το νευρωνικό δίκτυο MLP απαιτεί τουλάχιστον 150--200 δείγματα για να ξεπεράσει το $R^2 > 0.80$, ενώ το Random Forest περιορίζεται στο $R^2 \approx 0.47$ λόγω της αδυναμίας των δέντρων απόφασης να παρεμβάλλουν ομαλά σε συνεχή πολυδιάστατα πεδία.

\subsection{Ολοκληρωμένη Σύγκριση Αρχιτεκτονικών Surrogate}
Για την επιλογή του βέλτιστου κινητήρα υποκατάστασης αξιολογήθηκαν τέσσερις διαφορετικές κατηγορίες αλγορίθμων μηχανικής μάθησης: Gaussian Process (Matérn 5/2 ARD), Τεχνητό Νευρωνικό Δίκτυο MLP ($64 \times 32$), Random Forest (100 estimators) και $k$-Πλησιέστεροι Γείτονες ($k$-NN, $k=5$).

Στο Σχήμα~\ref{fig:surrogate_model_comparison} παρουσιάζεται η συγκριτική αξιολόγηση 6 διαστάσεων (ποιότητα προσαρμογής, μετρικές σφάλματος, κατανομή υπολοίπων, καθυστέρηση και ρυθμαπόδοση), ενώ στον Πίνακα~\ref{tab:surrogate_benchmark} καταγράφονται τα ποσοτικά αποτελέσματα.

\begin{figure}[htbp]
    \centering
    \includegraphics[width=0.98\textwidth]{figures/fig2_surrogate_model_comparison.png}
    \caption{Ολοκληρωμένη σύγκριση τεσσάρων αρχιτεκτονικών surrogate (GP, MLP, Random Forest, $k$-NN): (α) Συντελεστής $R^2$, (β) RMSE, (γ) MAE, (δ) Καθυστέρηση ανά 1.000 προβλέψεις, (ε) Ρυθμαπόδοση (eval/s) και (στ) Κατανομή υπολοίπων.}
    \label{fig:surrogate_model_comparison}
\end{figure}

../../evaluation/figures/table9_surrogate_models_and_gpu_benchmarks.tex

Από τον Πίνακα~\ref{tab:surrogate_benchmark} αναδεικνύεται η σαφής υπεροχή του Gaussian Process στην ακρίβεια ($R^2 = 0.8790$, $\text{RMSE} = 0.7396$), καθώς και η παροχή εγγενούς εκτίμησης αβεβαιότητας $\sigma(\mathbf{x})$, γεγονός κρίσιμο για την ασφαλή προώθηση στο Tier 3. Από την άλλη πλευρά, το μοντέλο MLP προσφέρει εξαιρετικά υψηλή ρυθμαπόδοση ($>1.2 \times 10^6\text{ eval/s}$ με καθυστέρηση μόλις $0.82\text{ ms}$ ανά 1.000 δείγματα), καθιστώντας το ιδανικό για μαζική παράλληλη αξιολόγηση υποψηφίων.

\subsection{Ανάλυση Ευαισθησίας Παραμέτρων μέσω Αυτόματου Προσδιορισμού Σχετικότητας (ARD)}
Η εκπαίδευση του GP με πυρήνα Matérn 5/2 και Αυτόματο Προσδιορισμό Σχετικότητας (Automatic Relevance Determination - ARD) επιτρέπει την εκμάθηση διακριτής κλίμακας μήκους (length-scale $\ell_d$) για κάθε μία από τις διαστάσεις του σχεδιαστικού χώρου. Μικρή κλίμακα μήκους $\ell_d$ σημαίνει ότι μικρή μεταβολή της παραμέτρου επιφέρει δραστική αλλαγή στην αεροδυναμική απόδοση, δηλαδή υψηλή σχετική ευαισθησία.

Στο Σχήμα~\ref{fig:surrogate_parameter_sensitivity} και στον Πίνακα~\ref{tab:gp_ard_sensitivity} αναλύεται η κατάταξη των σχεδιαστικών παραμέτρων ως προς την αεροδυναμική τους ευαισθησία.

\begin{figure}[htbp]
    \centering
    \includegraphics[width=0.95\textwidth]{figures/fig2_surrogate_parameter_sensitivity.png}
    \caption{Ανάλυση ευαισθησίας σχεδιαστικών παραμέτρων μέσω ARD GP: Κατάταξη παραμέτρων με βάση τις εκμαθημένες κλίμακες μήκους $\ell_d$ και αποτύπωση των επιμέρους μερικών επιδράσεων στην αεροδυναμική απόδοση.}
    \label{fig:surrogate_parameter_sensitivity}
\end{figure}

../../evaluation/figures/table_fig2_surrogate_parameter_sensitivity.tex

Τα ευρήματα του Πίνακα~\ref{tab:gp_ard_sensitivity} συμφωνούν απόλυτα με τις αρχές της αεροδυναμικής θεωρίας: η μέγιστη καμπυλότητα της αεροτομής (\texttt{naca\_m}, $40.3\%$), η γωνία συστροφής (\texttt{wing\_twist}, $16.2\%$), το πάχος της αεροτομής (\texttt{naca\_t}, $16.2\%$) και η διάμετρος της ατράκτου (\texttt{fuse\_max\_diam}, $12.9\%$) κυριαρχούν στην επίδοση, καθώς επηρεάζουν άμεσα την άντωση, την αποκόλληση της ροής και την παρασιτική οπισθέλκουσα. Αντίθετα, παράμετροι όπως το μήκος της ατράκτου και το εκπέτασμα (σε καθεστώς σταθερής επιφάνειας) παρουσιάζουν πολύ μεγαλύτερα length-scales ($\ell > 60$), συνεισφέροντας λιγότερο από $0.1\%$ στην τοπική διακύμανση.

\subsection{Υπολογιστική Επιτάχυνση Υλικού: CPU OpenMP έναντι NVIDIA CUDA GPU}
Για την επιτάχυνση των υπολογισμών γραμμικής άλγεβρας του Gaussian Process (κατασκευή πίνακα συνδιακύμανσης $K$ διαστάσεων $N \times N$ και επίλυση γραμμικού συστήματος Cholesky $K^{-1}\mathbf{y}$), συγκρίθηκε η εκτέλεση σε πολυνηματικό επεξεργαστή μέσω OpenMP (12 νήματα) έναντι υλοποίησης σε κάρτα γραφικών NVIDIA CUDA.

Στο Σχήμα~\ref{fig:surrogate_cuda_vs_cpu} και στο κάτω μέρος του Πίνακα~\ref{tab:surrogate_benchmark} συγκρίνεται ο χρόνος εκτέλεσης σε συνάρτηση με το μέγεθος δείγματος $N$.

\begin{figure}[htbp]
    \centering
    \includegraphics[width=0.92\textwidth]{figures/fig2_surrogate_cuda_vs_cpu.png}
    \caption{Σύγκριση υπολογιστικής επίδοσης επιταχυντή: Πολυνηματική CPU (OpenMP, 12 νήματα) έναντι NVIDIA CUDA GPU για τον υπολογισμό πίνακα συνδιακύμανσης $D=11$ και επίλυση συστήματος σε συνάρτηση με το μέγεθος $N$.}
    \label{fig:surrogate_cuda_vs_cpu}
\end{figure}

Για μικρά μεγέθη δειγμάτων ($N \le 360$), η εκτέλεση στην CPU υπερτερεί ελαφρώς ($0.14\text{ ms}$ έναντι $0.30\text{ ms}$ στο $N=100$), καθώς η επιβάρυνση εκκίνησης των πυρήνων CUDA (kernel launch overhead) και η μεταφορά μνήμης PCIe κυριαρχούν. Ωστόσο, όταν το μέγεθος κλιμακώνεται σε $N \ge 1.000$ δείγματα, η GPU αναλαμβάνει πλήρως την υπολογιστική υπεροχή ($6.43\text{ ms}$ έναντι $7.26\text{ ms}$, επιτάχυνση $1.13\times$), επιβεβαιώνοντας τη χρησιμότητά της σε σενάρια μαζικής επιγραμμικής επανεκπαίδευσης.
